\documentclass[12pt]{article}
\usepackage[margin=1in]{geometry}
\usepackage{amsmath, amssymb, amsfonts}
\usepackage{enumitem}
\usepackage{booktabs}
\usepackage{hyperref}
\hypersetup{colorlinks=true, linkcolor=blue, urlcolor=blue}

\setlist[enumerate]{leftmargin=*}

\title{\textbf{Quiz 1} \vspace{0.4cm}\\ECON 101: Macroeconomics}
\author{Instructor: Muhebullah Karimzada}
\date{September 28, 2026}

\begin{document}
\maketitle




\begin{enumerate}

    \item The fundamental economic problem arises because:
    \begin{enumerate}[label=\Alph*.]
        \item People have unlimited resources but limited wants
        \item Society must make choices because resources are scarce relative to wants
        \item Governments control the allocation of all resources
        \item Firms cannot produce enough goods to satisfy every consumer
    \end{enumerate}

    \item Which of the following is a \textbf{positive economic statement}?
    \begin{enumerate}[label=\Alph*.]
        \item The federal government should reduce income taxes.
        \item Canada should increase spending on education.
        \item An increase in the price of gasoline reduces the quantity of gasoline demanded, ceteris paribus.
        \item The government ought to provide more affordable housing.
    \end{enumerate}

    \item Opportunity cost is best described as:
    \begin{enumerate}[label=\Alph*.]
        \item The total monetary cost of producing a good
        \item The value of the next-best alternative that must be given up
        \item The amount of money spent on a decision
        \item The difference between total revenue and total cost
    \end{enumerate}

    \item A country's PPF is bowed outward from the origin primarily because:
    \begin{enumerate}[label=\Alph*.]
        \item All resources are equally suited to producing both goods
        \item The opportunity cost of producing additional units of one good tends to rise
        \item Resources become unlimited as production increases
        \item The economy experiences no scarcity
    \end{enumerate}

    \item Suppose Country X can produce either 40 units of wheat or 80 units of clothing, while Country Y can produce either 60 units of wheat or 90 units of clothing. Mutually beneficial trade is possible when countries specialize according to:
    \begin{enumerate}[label=\Alph*.]
        \item Absolute advantage only
        \item Comparative advantage
        \item Equal production possibilities
        \item Government regulation of trade
    \end{enumerate}

\end{enumerate}




\begin{enumerate}[resume]

    \item A student has 12 hours available for either studying or working. The student earns \$20 per hour at a part-time job. If the student chooses to study for one additional hour instead of working, what is the \textbf{opportunity cost} of that additional hour of studying? \textbf{(2 points)}

    \item Consider the following simplified PPF data:

    \begin{center}
    \begin{tabular}{ccc}
    \toprule
    Choice & Bicycles & Tablets \\
    \midrule
    A & 0   & 120 \\
    B & 20  & 100 \\
    C & 40  & 80  \\
    D & 60  & 60  \\
    E & 80  & 40  \\
    F & 100 & 20  \\
    \bottomrule
    \end{tabular}
    \end{center}

    \begin{enumerate}
        \item What is the opportunity cost of increasing bicycle production from 40 to 60? \textbf{(2 points)}
        \item Based on the table, are the opportunity costs constant, increasing, or decreasing? Briefly explain. \textbf{(2 points)}
    \end{enumerate}

    \item Country A can produce a maximum of 80 tons of wheat or 40 tons of steel. Country B can produce a maximum of 90 tons of wheat or 60 tons of steel.

    \begin{enumerate}
        \item Which country has the \textbf{absolute advantage} in wheat and which country has the absolute advantage in steel? \textbf{(2 points)}
        \item Calculate the opportunity cost of producing one ton of steel in each country. Which country has the \textbf{comparative advantage} in steel? \textbf{(2 points)}
    \end{enumerate}

\end{enumerate}

\end{document}
